% Exam Template for UMTYMP and Math Department courses
%
% Using Philip Hirschhorn's exam.cls: http://www-math.mit.edu/~psh/#ExamCls
%
% run pdflatex on a finished exam at least three times to do the grading table on front page.
%
%%%%%%%%%%%%%%%%%%%%%%%%%%%%%%%%%%%%%%%%%%%%%%%%%%%%%%%%%%%%%%%%%%%%%%%%%%%%%%%%%%%%%%%%%%%%%%

% These lines can probably stay unchanged, although you can remove the last
% two packages if you're not making pictures with tikz.
\documentclass[11pt]{exam}
\RequirePackage{amssymb, amsfonts, amsmath, latexsym, verbatim, xspace, setspace}
\RequirePackage{tikz, pgflibraryplotmarks}

% By default LaTeX uses large margins.  This doesn't work well on exams; problems
% end up in the "middle" of the page, reducing the amount of space for students
% to work on them.
\usepackage[margin=1in]{geometry}
\usepackage{enumerate}
\usepackage{amsthm}

\theoremstyle{definition}
\newtheorem{soln}{Solution}

% Here's where you edit the Class, Exam, Date, etc.
\newcommand{\class}{Math 350 Section 1}
\newcommand{\term}{Fall 2024}
\newcommand{\examnum}{Final Exam}
\newcommand{\examdate}{December 16, 2024}
\newcommand{\timelimit}{1 Hour 50 Minutes}
\newcommand{\ol}[1]{\overline{#1}}

% For an exam, single spacing is most appropriate
\singlespacing
% \onehalfspacing
% \doublespacing

% For an exam, we generally want to turn off paragraph indentation
\parindent 0ex

\begin{document} 

% These commands set up the running header on the top of the exam pages
\pagestyle{head}
\firstpageheader{}{}{}
\runningheader{\class}{\examnum\ - Page \thepage\ of \numpages}{\examdate}
\runningheadrule

\begin{flushright}
\begin{tabular}{p{2.8in} r l}
\textbf{\class} & \textbf{Name (Print):} & \makebox[2in]{\hrulefill}\\
\textbf{\term} &&\\
\textbf{\examnum} & \textbf{Student ID:}&\makebox[2in]{\hrulefill}\\
\textbf{\examdate} &&\\
\textbf{Time Limit: \timelimit} % & Teaching Assistant & \makebox[2in]{\hrulefill}
\end{tabular}\\
\end{flushright}
\rule[1ex]{\textwidth}{.1pt}


This exam contains \numpages\ pages (including this cover page) and
\numquestions\ problems.  Check to see if any pages are missing.  Enter
all requested information on the top of this page, and put your initials
on the top of every page, in case the pages become separated.\\

You may \textit{not} use your books or notes on this exam.

You are required to show your work on each problem on this exam.  The following rules apply:\\

\begin{minipage}[t]{3.7in}
\vspace{0pt}
\begin{itemize}

%\item \textbf{If you use a ``fundamental theorem'' you must indicate this} and explain
%why the theorem may be applied.

\item \textbf{Organize your work}, in a reasonably neat and coherent way, in
the space provided. Work scattered all over the page without a clear ordering will 
receive very little credit.  

\item \textbf{Mysterious or unsupported answers will not receive full
credit}.  A correct answer, unsupported by calculations, explanation,
or algebraic work will receive no credit; an incorrect answer supported
by substantially correct calculations and explanations might still receive
partial credit.  This especially applies to limit calculations.

\item If you need more space, use the back of the pages; clearly indicate when you have done this.

\item \textbf{Box Your Answer} where appropriate, in order to clearly indicate what you consider the answer to the question to be.

\item \textbf{Default metric}: unless otherwise stated, the set $\mathbb{R}$ or $\mathbb{R}^n$ will be considered a metric space with the Euclidean metric
\end{itemize}

Do not write in the table to the right.
\end{minipage}
\hfill
\begin{minipage}[t]{2.3in}
\vspace{0pt}
%\cellwidth{3em}
\gradetablestretch{2}
\vqword{Problem}
\addpoints % required here by exam.cls, even though questions haven't started yet.	
\gradetable[v]%[pages]  % Use [pages] to have grading table by page instead of question

\end{minipage}
\newpage % End of cover page

%%%%%%%%%%%%%%%%%%%%%%%%%%%%%%%%%%%%%%%%%%%%%%%%%%%%%%%%%%%%%%%%%%%%%%%%%%%%%%%%%%%%%
%
% See http://www-math.mit.edu/~psh/#ExamCls for full documentation, but the questions
% below give an idea of how to write questions [with parts] and have the points
% tracked automatically on the cover page.
%

%%%%%%%%%%%%%%%%%%%%%%%%%%%%%%%%%%%%%%%%%%%%%%%%%%%%%%%%%%%%%%%%%%%%%%%%%%%%%%%%%%%%%

\begin{questions}

\addpoints

\question[10]\mbox{}
\textbf{TRUE or FALSE!}  Write  TRUE if the statement is true.  Otherwise, write FALSE.  Your response should be in ALL CAPS.  No justification is required.
\begin{enumerate}[(a)]
\item Every subset of $\mathbb{R}$ with the discrete metric is open
\vspace{1.2in}
\item The function $f(x) = e^x+x$ has an $x$-intercept
\vspace{1.2in}
\item If $f(x)$ is differentiable at $x=0$, then so is $f(x^2)$
\vspace{1.2in}
\item The set of all irrational numbers is uncountable
\vspace{1.2in}
\item In any metric space, compact sets must be closed
\vspace{1.2in}
\end{enumerate}

\newpage
\question[10]\mbox{}
Let $(M,d)$ be a metric space and $K\subseteq M$.
\begin{enumerate}[(a)]
\item Write the definition of $K$ being compact
\vspace{2in}
\item Write the definition of $K$ being bounded
\vspace{2in}
\item Prove that if $K$ is not bounded, then it cannot be compact
\end{enumerate}

\newpage
\question[10]\mbox{}
In this problem, let $\alpha(x) = x$ and
$$f(x) = \left\lbrace\begin{array}{cc}
1, & x\neq 3\\
5, & x =   3
\end{array}\right.$$

Carefully prove $f$ is integrable with respect to $\alpha$ on $[0,3]$ and determine the value of $\int_0^3 fd\alpha$

\newpage
\question[10]\mbox{}

Consider the function
$$f(x) = \left\lbrace\begin{array}{cc}
x, & x\in \mathbb{Q}\\
0, & x\notin \mathbb{Q}
\end{array}\right.$$

\begin{enumerate}[(a)]
\item
Prove that $f(x)$ is not continuous at $x=c$ for $c\neq 0$
\vspace{3in}
\item
Prove that $f(x)$ is continuous at $x=0$
\end{enumerate}

\newpage
\question[10]\mbox{}

\begin{enumerate}[(a)]
\item
Carefully explain why there is no continuous, surjective function $f: [0,1]\rightarrow (0,1]$
\vspace{2.5in}
\item
Carefully explain why there is no continuous, surjective function $f: [0,1]\rightarrow [0,1/2)\cup (1/2,1]$
\vspace{2.5in}
\item
Show that there is a continuous, surjective function $f: (0,1]\rightarrow [0,1]$
\end{enumerate}


\newpage
\question[10]\mbox{}
Let $\{x_n\}$ be a monotone increasing sequence of real numbers.

\begin{enumerate}[(a)]
\item
State the Completeness Axiom
\vspace{2in}
\item
Prove that if $\{x_n\}$ is also bounded above, then it converges.
\end{enumerate}



\newpage
\question[10]\mbox{}
Suppose that $f(x)$ is a polynomial with roots at $x=1$, $x=2$, and $x=3$

\begin{enumerate}[(a)]
\item
State the Mean Value Theorem
\vspace{2in}
\item
Prove that $f'(x)$ has at least two different real roots
\vspace{3.5in}
\item
Prove that $f''(x)$ has at least one real root
\end{enumerate}


\end{questions}

\end{document}

